%% ==========================================================================
%%  SECTION 8 --- THE GAP AND OUR DIRECTION                     [Track T4]
%%
%%  Job: everything before this has been setup. Name the gap, state the two
%%  concrete research questions arising from it, and say what we will build.
%%  ~2 pages. This section is what the survey is FOR, and it is the one the
%%  instructor will read most closely.
%%
%%  SOURCE: report/notes/_synthesis.md holds the argument in draft, including
%%  the list of things that would KILL it. Do not write this section before
%%  the synthesis meeting on 23 August.
%% ==========================================================================
\section{The Gap and Our Direction}
\label{sec:gap}

\subsection{The gap}
\label{sec:gap:statement}

\todo{One paragraph, stated flatly, drawing \S4, \S5 and \S6 together:
      private training and private retrieval are two mature literatures that
      have developed in parallel; the only system spanning both predates the
      whole of the second; and the current best training system explicitly
      declines the fetch problem. Nobody has plugged modern private-retrieval
      machinery into a privately trained recommender.
      Then the honest qualifier --- this is a gap in the ENGINEERING and in
      the ANALYSIS, not an open theory problem. Claiming otherwise would be
      the easiest way to lose credibility here.}

\begin{figure}[t]
\centering
\todo{Figure: the two literature timelines side by side, 2013--2026, with the
      single arrow between them (\pirsonaplain{}, 2021) and the dangling
      ``other means'' arrow from \nudgeplain{}. This figure IS the argument;
      if it works, the section could almost be read from it alone. TikZ.}
\caption{Two parallel literatures, and the gap between them.}
\label{fig:gap:timelines}
\end{figure}

\subsection{Research question 1: where is the crossover?}
\label{sec:gap:n1}

\todo{Develop the observation planted in \Cref{sec:retrieval:asymmetry}.
      \nudgeplain{} publishes $\itemembed \in \R^{\dime \times \items}$ in the
      clear, and user $i$ holds its own ratings $\uvec{i}$. So the user can
      download $\itemembed$ once and compute its own top-$\topk$ LOCALLY, with
      no cryptography. At MovieLens scale ($\dime=20$, $\items\approx 4000$)
      that is under a megabyte.
      But $\itemembed$ grows with the catalogue, and at Criteo scale
      (hundreds of millions of items) downloading it is absurd --- at which
      point the \S5 machinery becomes necessary.
      THE QUESTION: as a function of $\items$, $\dime$, and client bandwidth,
      where do the boundaries lie between (a) download $\itemembed$ and
      compute locally, (b) private ANN over $\itemembed$, and (c) PIR fetch
      against the catalogue?
      Say why this is worth answering: every system in \S5 assumes its own
      regime, and nobody has drawn the boundary. It is also cheap --- mostly
      analysis plus measurement. \verify{confirm no prior work draws this
      boundary before claiming novelty}}

\subsection{Research question 2: what does a public $\itemembed$ leak?}
\label{sec:gap:n2}

\todo{Develop the thread from \Cref{sec:threat:leakage}. Every server holds
      $\itemembed$ in the clear --- a complete latent-factor model of the
      catalogue. So an observed fetch is not one bit about a user; it is a
      PROJECTION ONTO A KNOWN BASIS, and a handful of fetches constrains the
      taste vector sharply.
      Make it concrete and falsifiable: given public $\itemembed$ and $j$
      observed fetches, reconstruct $\hat{\avec{i}}$ by least squares and
      measure $\cos(\hat{\avec{i}}, \avec{i})$ as a function of $j$.
      Then the payoff: this is the QUANTITATIVE argument for why the delivery
      layer must be private. It turns our contribution from an engineering
      preference into a demonstrated necessity.
      Note also that \nudgeplain{} \S9 offers DP on $\itemembed$ as a
      mitigation, so the natural follow-up is what that costs in nDCG and
      whether it actually blunts the attack.}

\subsection{What we propose to build}
\label{sec:gap:proposal}

\todo{Keep this SHORT and do not over-commit --- the survey is not the design
      document, and the architecture is settled at the instructor meeting.
      Say: a private delivery layer for a privately-trained recommender,
      composing a \nudgeplain{}-style training core with a DPF-PIR fetch layer
      and \pirsonaplain{}'s harvest loop, evaluated end to end under the
      standard set in \Cref{sec:eval:ourstandard}.
      Then state both open decisions HONESTLY rather than pretending they are
      settled:
      (a) whether to build on the \nudgeplain{} reference artifact
          (\url{https://github.com/NudgeArtifact/private-recs}, MIT-licensed
          Go) or reimplement the training core;
      (b) whether the serving layer needs cryptography at all in our
          parameter regime --- which is exactly \Cref{sec:gap:n1}.
      Presenting both framings neutrally is the point; the instructor meeting
      decides.}

\subsection{Threats to this reading of the literature}
\label{sec:gap:threats}

\todo{Include this subsection. A survey that only looks for confirming
      evidence is worthless, and saying so explicitly is the strongest
      credibility signal available in a document with no experimental results.
      List honestly:
      (1) someone may already have composed a private-ANN backend with a
          private-MF recommender --- \Cref{sec:both:empty} depends on a
          forward-citation sweep we must actually run;
      (2) the \Cref{sec:gap:n1} crossover may sit so far out that
          $\itemembed$ is downloadable at every realistic catalogue size,
          reducing RQ1 to a one-line observation;
      (3) the \Cref{sec:gap:n2} reconstruction may already be folded into the
          standard embedding-inversion literature.
      For each, say what evidence would settle it.}
